\documentclass[10pt,a4paper,twocolumn,twoside]{article}
\usepackage[margin=0.75in,columnsep=0.3in]{geometry}
\usepackage[utf8]{inputenc}
\usepackage[T1]{fontenc}
\usepackage{amsmath,amssymb,graphicx,booktabs,hyperref,xcolor}
\newcommand{\todo}[1]{\textcolor{red}{[TODO: #1]}}

\title{Mass-Effect Deformation and Hierarchy-Aware Loss for Label-Free\\
Synthetic-Tumor Brain MRI Segmentation}
\author{Piyush Garg, Sahaj Sharma, Goutam Jain\\
\small Supervisor: Dr.\ Vijay\\
\small Netaji Subhas University of Technology (NSUT), New Delhi}
\date{}

\begin{document}
\maketitle

\begin{abstract}
Training deep networks to segment brain tumors normally needs hundreds of
voxel-wise annotated scans, which are slow and costly to produce.
Label-free approaches avoid this by training only on synthetic data:
fake tumors are pasted into healthy scans, and the network never sees a
human-labeled tumor. LF-SynthSeg \cite{lfsynthseg} follows this idea for
T2-weighted MRI by filling random ellipsoid-based tumor shapes with the
intensities of healthy brain tissue and training a 3D network on the result.
We re-implement this pipeline from scratch (healthy HCP scans, tissue
labels from FSL FAST, a bank of 1000 synthetic tumors, a DynUNet
segmenter) and evaluate it on 31 BraTS~2020 cases using the Dice score,
HD95, AUPRC and sensitivity. On top of this baseline we study two
additions: (i) a 3D mass-effect module that pushes surrounding tissue away
from the synthetic tumor with a smooth, folding-free deformation, and
(ii) a hierarchy-aware loss that exploits the nested structure of tumor
sub-regions (necrotic core, enhancing tumor, edema). \todo{one sentence on
the reproduced baseline result, e.g.\ ``Our reproduction reaches a Dice of X
against 0.823 reported originally''.} \todo{one sentence per addition with
its effect and significance over 3+ seeds, or an honest negative result.}
\end{abstract}

\section{Introduction}
Brain tumors such as gliomas are among the most lethal cancers, and magnetic
resonance imaging (MRI) is the main tool for diagnosing them and planning
treatment. Delineating the tumor on each scan is a time-consuming task that
varies between observers, so automatic segmentation is attractive.
Convolutional networks, in particular U-Net variants
\cite{unet,nnunet}, now reach expert-level accuracy on public benchmarks such as
the Brain Tumor Segmentation (BraTS) challenge \cite{menze,bakas}.

These models, however, rely on large sets of manually annotated scans.
Annotations need trained radiologists, and a model trained on one scanner,
protocol or hospital often loses accuracy elsewhere. Synthetic training data
offers a way out. SynthSeg \cite{synthseg} showed that a network trained
only on synthetic scans generated from label maps can segment healthy brain
anatomy across contrasts and resolutions. For tumors, label-free
tumor synthesis was explored for liver CT \cite{labelfree_liver}, and
LF-SynthSeg \cite{lfsynthseg} brought it to brain MRI: random tumor shapes are
filled with the intensities of CSF, gray matter and white matter taken from
the same scan, pasted into healthy T2 images, and a segmentation network is
trained on these images with no real tumor annotation at all.

Two aspects of this approach are open for improvement. First, a real
tumor does not only replace healthy tissue; it also \emph{displaces} it. This
\emph{mass effect} pushes the ventricles and cortex away from the lesion, and
the synthetic tumors of the original method do not model it. Second, tumors
have a nested structure (a whole tumor contains a core, which contains
enhancing and necrotic parts), yet the loss used for training treats
the classes as unrelated.

In this work we make the following contributions:
\begin{itemize}
\item We re-implement LF-SynthSeg from scratch on public data (HCP T2 scans
as healthy hosts, FSL FAST tissue labels, a bank of 1000 synthetic tumors, a
3D DynUNet) and report its performance on 31 BraTS~2020 cases under the
original protocol.
\item We add a 3D mass-effect deformation that moves surrounding tissue away
from the synthetic tumor while keeping the transformation smooth and
invertible (positive Jacobian determinant) and confined to the brain.
\item We add a hierarchy-aware loss that uses the nesting of tumor
sub-regions, and compare variants in a controlled ablation with several
random seeds and paired statistical tests.
\item We document the practical problems met while reproducing the method,
such as the gap between synthetic and real images, and how we addressed
them without using the test cases for tuning.
\end{itemize}

\section{Literature Review}
\subsection{Supervised brain tumor segmentation}
The BraTS challenge provides multi-institutional multi-parametric scans with
expert labels for tumor sub-regions \cite{menze,bakas}. Early deep
approaches used patch-wise 2D and 3D convolutional networks; the encoder-decoder
U-Net \cite{unet} and its 3D extensions became the standard. The nnU-Net
framework \cite{nnunet} showed that careful automatic configuration of a
plain U-Net is hard to beat, and winners of BraTS used region-based, nested
targets (whole tumor, tumor core, enhancing tumor) together with Dice-based
losses \cite{myronenko,vnet}. All of these need the full set of manual
labels.

\subsection{Synthetic data and data augmentation for lesions}
Generative adversarial networks have been used to create additional training
images \cite{shin}, but they need real data to train and can copy its biases.
Mixing-based augmentation such as CarveMix \cite{carvemix} combines lesions
from labeled scans with other images, which still requires some annotated
lesions. Domain randomization takes a different route: SynthSeg
\cite{synthseg} generates scans with random contrast, resolution and
artifacts from label maps, and a network trained purely on them generalizes
to real scans. This idea is the base of label-free methods.

\subsection{Label-free tumor synthesis}
Hu et al.\ \cite{labelfree_liver} trained liver tumor segmentation on tumors
synthesized in healthy CT scans, with shapes and textures designed by hand and
no real tumor labels. LF-SynthSeg \cite{lfsynthseg} applies the principle to
T2-weighted brain MRI. Random shapes are built as unions of ellipsoids, their
inner regions are linked to the three healthy tissues (necrotic core with CSF,
enhancing tumor with gray matter, edema with white matter), missing
intensities are completed by interpolation, and the result is blended into
healthy scans whose tissue labels come from an automatic segmentation. The
authors report a whole-tumor Dice of 0.823 on BraTS~2020. Our work is a
from-scratch reproduction and an extension of this method.

\subsection{Mass effect and hierarchy-aware learning}
Biophysical models of tumor growth simulate the deformation of brain tissue
around a lesion \cite{hogea}, but are expensive and are usually not used for
data synthesis at scale. A lighter approach is a smooth displacement field
that moves tissue radially away from the tumor boundary; we explore this
for synthetic data. For the loss, nested label structure has been exploited in
BraTS through region-based targets \cite{nnunet,myronenko}; we test such
ideas in the label-free setting.

\subsection{Gap addressed by this work}
To our knowledge, the effect of explicit mass effect and of a hierarchy-aware
loss on label-free brain tumor synthesis has not been studied.


\begin{thebibliography}{99}\small
\bibitem{lfsynthseg} Y.~Xu \emph{et al.}, ``LF-SynthSeg: Label-free brain tissue-assisted tumor synthesis and segmentation,'' \emph{IEEE Journal of Biomedical and Health Informatics}, 2025.
\bibitem{synthseg} B.~Billot \emph{et al.}, ``SynthSeg: Segmentation of brain MRI scans of any contrast and resolution without retraining,'' \emph{Medical Image Analysis}, vol.~86, 2023.
\bibitem{labelfree_liver} Q.~Hu \emph{et al.}, ``Label-free liver tumor segmentation,'' in \emph{Proc.\ IEEE/CVF CVPR}, 2023.
\bibitem{menze} B.~H. Menze \emph{et al.}, ``The multimodal brain tumor image segmentation benchmark (BRATS),'' \emph{IEEE Trans.\ Medical Imaging}, vol.~34, no.~10, 2015.
\bibitem{bakas} S.~Bakas \emph{et al.}, ``Advancing The Cancer Genome Atlas glioma MRI collections with expert segmentation labels and radiomic features,'' \emph{Scientific Data}, vol.~4, 2017.
\bibitem{unet} O.~Ronneberger, P.~Fischer, T.~Brox, ``U-Net: Convolutional networks for biomedical image segmentation,'' in \emph{MICCAI}, 2015.
\bibitem{nnunet} F.~Isensee \emph{et al.}, ``nnU-Net: a self-configuring method for deep learning-based biomedical image segmentation,'' \emph{Nature Methods}, vol.~18, 2021.
\bibitem{myronenko} A.~Myronenko, ``3D MRI brain tumor segmentation using autoencoder regularization,'' in \emph{BrainLes (MICCAI Workshop)}, 2018.
\bibitem{vnet} F.~Milletari, N.~Navab, S.-A. Ahmadi, ``V-Net: Fully convolutional neural networks for volumetric medical image segmentation,'' in \emph{3DV}, 2016.
\bibitem{shin} H.-C. Shin \emph{et al.}, ``Medical image synthesis for data augmentation and anonymization using generative adversarial networks,'' in \emph{SASHIMI (MICCAI Workshop)}, 2018.
\bibitem{carvemix} X.~Zhang \emph{et al.}, ``CarveMix: A simple data augmentation method for brain lesion segmentation,'' \emph{NeuroImage}, vol.~271, 2023.
\bibitem{hogea} C.~Hogea, C.~Davatzikos, G.~Biros, ``An image-driven parameter estimation problem for a reaction--diffusion glioma growth model with mass effects,'' \emph{J.\ Mathematical Biology}, vol.~56, 2008.
\end{thebibliography}
\end{document}
